O crescente volume de dados em aplicações web, especialmente no âmbito da Administração Pública, demanda o desenvolvimento de interfaces que facilitem a sua organização, visualização e interpretação. Nesse contexto, os Sistemas de Informação Geográfica (SIG) emergem como ferramentas cruciais para a gestão urbana e a transparência de dados públicos. No entanto, o desenvolvimento de suas interfaces apresenta desafios únicos relacionados ao desempenho no carregamento de grandes volumes de dados espaciais, à responsividade em múltiplos dispositivos e à garantia de uma experiência de usuário (UX) intuitiva. Este trabalho se insere no projeto "Sobral em Mapas" e aborda a construção de sua plataforma web. O objetivo geral é desenvolver um sistema interativo e responsivo para a produção, edição e visualização de mapas digitais associados a dados abertos da Prefeitura Municipal de Sobral. A aplicação foi inteiramente estruturada sobre o framework PHP Laravel, que gerencia de forma integrada tanto as funcionalidades do back-end quanto a apresentação do front-end através do padrão arquitetural Model-View-Controller (MVC). Para a interface, foram utilizadas tecnologias como HTML, CSS e JavaScript, complementadas pelo framework Bootstrap. Como principal contribuição, este trabalho apresenta a implementação da arquitetura central da plataforma, estabelecendo uma base tecnológica robusta. O resultado é um sistema funcional que viabiliza a visualização de dados geoespaciais de forma clara e eficiente, garantindo uma experiência de usuário acessível e, fundamentalmente, servindo como um alicerce sólido para futuras implementações e melhorias no âmbito do projeto "Sobral em Mapas".

% Separe as palavras-chave por ponto
\palavraschave{Desenvolvimento Web, Mapas Digitais, Sistemas de Informação Geográfica, Laravel, Gestão Pública, Dados Abertos.}